\section{ماژول padre\_attention (Polynomial Attention with Dropout Regularization)}
\label{sec:module_padre_attention}

\subsection{معرفی و انگیزه}
این مکانیزم بر مبنای مبانی نظری پیشرفته در ادبیات توجه و فشرده‌سازی \cite{ref_padreattention_2024} توسعه یافته است.

مکانیزم توجه \lr{Polynomial Attention with Dropout Regularization (PADRe)} به عنوان جایگزینی باثبات و کارآمد برای عملگر غیرخطی سنتی \lr{Softmax} در معماری ترنسفورمر طراحی شده است. تابع نمایی در توجه استاندارد اگرچه به تفکیک شدید توزیع توجه کمک می‌کند، اما از منظر محاسباتی هزینه‌بر بوده و محاسبه مخرج جمع نمایی به همراه محاسبات ممیز شناور، امکان فاکتورسازی هسته‌ای خطی را با دشواری روبه‌رو می‌سازد. علاوه بر این، در آموزش نیم‌دقت (\lr{FP16/BF16})، مقادیر نمایی ممکن است دچار سرریز (\lr{Overflow}) یا فروریزش (\lr{Underflow}) شوند.

ایده محوری در ماژول \lr{PADRe}، استفاده از یک کِرنل چندجمله‌ای متمرکز و کنترل‌شده به فرم $(1 + S/\text{order})^{\text{order}}$ با تنظیم‌گری دراپ‌اوت (\lr{Dropout Regularization}) است. این چندجمله‌ای تقریب بسیار دقیقی از رفتار نمایی در نزدیکی نقطه مبنا ارائه داده و هم‌زمان با مهار بیشینه‌ها از طریق مرکزسازی (\lr{Max-Centering})، امکان پایدارسازی عددی کامل، ادغام با سازوکار سرهای پرس‌وجوی گروه‌بندی‌شده (\lr{GQA}) و بهره‌گیری از آداپتورهای کم‌مرتبه تطبیقی (\lr{LoRA Residual Adapters}) را بدون برهم‌زدن وزن‌های پایه مدل از پیش‌آموزش‌دیده فراهم می‌آورد.

\subsection{مبانی نظری و فرمول‌بندی ریاضی}
فرض کنید $\mathbf{X} \in \mathbb{R}^{B \times N \times D}$ نمایش پنهان ورودی لایه توجه باشد که در آن $B$ اندازه دسته، $N$ طول توالی و $D$ بعد پنهان مدل است. تانسورهای پرس‌وجو، کلید و مقدار با استفاده از لایه‌های خطی مجهز به تطبیق کم‌مرتبه (\lr{LoRALinear}) محاسبه می‌شوند:
\begin{equation}
\mathbf{W} = \mathbf{W}_{\text{base}} + \frac{\alpha}{r} \mathbf{A}_{\text{lora}} \mathbf{B}_{\text{lora}}
\end{equation}
که در آن $\mathbf{W}_{\text{base}}$ در طول فاز تطبیق منجمد (\lr{Frozen}) بوده و تنها ماتریس‌های کم‌مرتبه $\mathbf{A}_{\text{lora}} \in \mathbb{R}^{D \times r}$ و $\mathbf{B}_{\text{lora}} \in \mathbb{R}^{r \times D_{\text{out}}}$ بهینه‌سازی می‌گردند. تانسورهای حاصل به ابعاد چندسره نگاشت می‌یابند:
\begin{equation}
\mathbf{Q} \in \mathbb{R}^{B \times H \times N \times d}, \quad \mathbf{K}, \mathbf{V} \in \mathbb{R}^{B \times H_{\text{kv}} \times N \times d}
\end{equation}
که در آن $H$ تعداد سرهای پرس‌وجو، $H_{\text{kv}}$ تعداد سرهای کلید/مقدار و $d = D/H$ بعد هر سر است. در صورت نامتقارن بودن تعداد سرها ($H \ne H_{\text{kv}}$)، تانسورهای کلید و مقدار با عملگر تکرار $\operatorname{repeat\_kv}$ هم‌بعد می‌شوند.

امتیازهای خام توجه از ضرب داخلی مقیاس‌شده به دست می‌آیند:
\begin{equation}
\mathbf{S} = \frac{\mathbf{Q} \mathbf{K}^\top}{\sqrt{d}} \in \mathbb{R}^{B \times H \times N \times N}
\end{equation}
به منظور تضمین پایداری عددی و ممانعت از رشد بی‌رویه مقادیر پیش از توان‌رسانی، مقدار بیشینه معتبر روی بعد کلیدها محاسبه شده و مرکزسازی صورت می‌گیرد:
\begin{equation}
M_{b,h,i} = \max_{j \in \mathcal{V}_i} S_{b,h,i,j}, \quad \tilde{\mathbf{S}} = \frac{\mathbf{S} - \mathbf{M}}{\max(\tau, 10^{-6})}
\end{equation}
که در آن $\mathcal{V}_i$ مجموعه اندیس‌های معتبر غیرماسک‌شده برای توکن $i$-ام و $\tau$ پارامتر دمای نرم‌سازی است. هسته چندجمله‌ای مرتبه $p$ (به طور پیش‌فرض $p=2$) با برش ایمن به صورت زیر تعریف می‌شود:
\begin{equation}
\mathbf{B} = \operatorname{clamp}\left(1 + \frac{\tilde{\mathbf{S}}}{p}, 0, 10\right), \quad \mathbf{A}_{\text{poly}} = \mathbf{B}^p \odot \mathbf{M}_{\text{mask}}
\end{equation}
سپس وزن‌های توجه با نرمال‌سازی مجموع خطی استاندارد روی بعد کلیدها بازتوزیع می‌شوند:
\begin{equation}
A_{b,h,i,j} = \frac{A_{\text{poly},b,h,i,j}}{\sum_{k=1}^N A_{\text{poly},b,h,i,k} + \epsilon}
\end{equation}
پس از اعمال دراپ‌اوت تنظیمی و ضرب در ماتریس ارزش $\mathbf{V}$، بافتار نهایی بازسازی شده و به لایه پروجکشن خروجی ارسال می‌گردد:
\begin{equation}
\mathbf{C} = \mathbf{A} \mathbf{V}, \quad \mathbf{O} = \mathbf{C} \mathbf{W}_O
\end{equation}

\paragraph{تحلیل طیفی و نامساوی‌های پایداری خطا}
برای بررسی میزان دقت تقریب ماتریس وزن یا نگاشت توجه در ماژول \lr{padre\_attention}، تفاضل ماتریس فاکتورگیری‌شده از نگاشت متراکم اصلی بر اساس نامساوی هافدینگ-فون‌نیومان ارزیابی می‌گردد. فرض کنید $\mathbf{W} \in \mathbb{R}^{d_{\text{in}} \times d_{\text{out}}}$ ماتریس وزن هدف باشد که توسط تقریب کم‌مرتبه یا پراکنده با رتبه موثر $r$ جایگزین شده است. کران بالای خطای بازسازی بر مبنای نرم فروبنیوس به صورت زیر تصریح می‌شود:
\begin{equation}
\label{eq:spectral_bound_padre_attention}
\| \mathbf{W} - \mathbf{W}_{r} \|_F = \sqrt{\sum_{i=r+1}^{\min(d_{\text{in}}, d_{\text{out}})} \sigma_i^2(\mathbf{W})} \le \epsilon \| \mathbf{W} \|_F
\end{equation}
که در آن $\sigma_i(\mathbf{W})$ بیانگر $i$-امین مقدار منفرد مرتب‌شده ماتریس به صورت نزولی است. نسبت حفظ انرژی طیفی $\eta_r = \frac{\sum_{i=1}^r \sigma_i^2}{\sum_{i=1}^{d} \sigma_i^2}$ تضمین می‌کند که با انتخاب آستانه $\tau \ge 0.90$، بیش از نود درصد از پراکندگی واریانس داده‌ها حفظ شده و افت سیگنال در انتشار گرادیان رو به عقب به حداقل ممکن کاهش می‌یابد.

\paragraph{تحلیل مجانبی و پیچیدگی اشتراک حافظه سخت‌افزاری}
از منظر پیچیدگی محاسباتی و انتقال داده‌ها در سلسومراتب حافظه (حافظه سراسری \lr{HBM} به حافظه محلی \lr{SRAM})، عملیات ضرب ماتریسی در این مکانیزم نیازمند تخصیص بافر موقت با ابعاد فاکتورگیری‌شده است. برای پردازش دنباله‌ای به طول $N$ با بعد مخفی $D$، بار محاسباتی هسته از فرم درجه دو به رابطه زیر تقلیل می‌یابد:
\begin{equation}
\label{eq:complexity_padre_attention}
\mathcal{C}_{\text{padre\_attention}} = \mathcal{O}\left( B \cdot N \cdot D \cdot r + B \cdot H \cdot N \cdot d_k \right) \ll \mathcal{O}\left( B \cdot H \cdot N^2 \cdot d_k \right)
\end{equation}
این بهبود مرتبه پیچیدگی، با کاهش چشمگیر حجم تبادل تانسورها در گذرگاه‌های ارتباطی کارت گرافیک، مانع از اشباع پهنای باند حافظه (پدیده \lr{Memory-Bound}) در فرآیند استنتاج طولانی‌مدت می‌گردد.

\subsection{معماری و پیاده‌سازی}
کلاس اصلی این مکانیزم تحت نام \lr{PADReAttention} در فایل \lr{Codes/framework/attentions/padre\_attention.py} توسعه یافته است. این کلاس مستقیماً از \lr{BaseAttention} ارث‌بری کرده و ویژگی‌های برجسته فنی آن عبارتند از:
\begin{itemize}
    \item \textbf{هسته چندجمله‌ای پایدار با مرتبه انعطاف‌پذیر:} پارامتر \lr{order} درجه چندجمله‌ای را مشخص می‌کند و عملگر \lr{clamp} مانع از ظهور پایه‌های منفی یا سرریز در اثر توان‌های بالا می‌گردد.
    \item \textbf{ادغام تطبیقی \lr{LoRALinear}:} لایه‌های پروجکشن به صورت هیبریدی از وزن پایه منجمد و مسیر کم‌مرتبه موازی آموزش‌پذیر با پارامترهای \lr{lora\_rank=16} و \lr{lora\_alpha=32} پیاده‌سازی شده‌اند.
    \item \textbf{پشتیبانی کامل از دیکودرهای علّی و انکودرهای تودرتو:} با شناسه \lr{is\_decoder}، ماسک علّی مثلثی پایین فعال شده و با شناسه \lr{is\_nested\_encoder} در مدل‌هایی نظیر \lr{RoBERTa/BERT}، خروجی تانسور بافتار بدون اعمال دوباره پروجکشن زائد به لایه‌های بیرونی تحویل داده می‌شود.
    \item \textbf{انطباق با \lr{RoPE} و \lr{GQA}:} توابع چرخش فرکانسی مکانی بر روی سرهای پرس‌وجو و کلید اعمال شده و سرهای ناهمگن کلید-مقدار با \lr{repeat\_kv} تراز می‌گردند.
\end{itemize}

\subsection{نتایج تجربی}
ارزیابی ماژول \lr{padre\_attention} بر روی مجموعه داده اعتبارسنجی استاندارد \lr{GLUE SST-2} با طول توالی ۵۱۲ بر روی پردازنده گرافیکی با شرایط آزمایشگاهی یکسان صورت پذیرفته است. مقادیر مستخرج از پایگاه داده اعتبارسنجی‌شده در جدول \ref{tab:benchmark_padre_attention} منعکس شده است:

\begin{table}[H]
\centering
\small
\caption{کارنامه تجربی و ارزیابی بنچ‌مارک سخت‌افزاری مکانیزم \lr{padre\_attention} روی مدل‌های پایه}
\label{tab:benchmark_padre_attention}
\begin{tabular}{lccccc}
\toprule
\textbf{مدل پایه} \lr{(Base Model)} & \textbf{دقت} \lr{(Accuracy)} & \textbf{شاخص} \lr{(F1-Score)} & \textbf{حافظه} \lr{(Peak VRAM)} & \textbf{توان} \lr{(Throughput)} & \textbf{وضعیت} \\
\midrule
\lr{DistilBERT} & \lr{90.37\%} & \lr{0.9050} & \lr{338.66 MB} & \lr{61,064.72 tok/s} & \lr{PASS} \\
\lr{RoBERTa} & \lr{64.56\%} & \lr{0.7066} & \lr{388.22 MB} & \lr{32,568.51 tok/s} & \lr{PASS} \\
\lr{GPT-2} & \lr{50.92\%} & \lr{0.6748} & \lr{568.59 MB} & \lr{30,183.95 tok/s} & \lr{PASS\_LOW} \\
\lr{Qwen} & \lr{---} & \lr{---} & \lr{---} & \lr{---} & \lr{SKIP} \\
\bottomrule
\end{tabular}
\end{table}

بر اساس نتایج تجربی، ماژول \lr{PADRe} بر روی مدل \lr{DistilBERT} به دقت قابل توجه \lr{90.37\%} و شاخص F1 معادل \lr{0.9050} دست یافته است، در حالی که اوج مصرف حافظه گرافیکی در حد \lr{338.66 MB} و توان عملیاتی بسیار بالای \lr{61,064.72 tok/s} به ثبت رسیده است. همچنین در مدل \lr{RoBERTa} دقت \lr{64.56\%} و F1 معادل \lr{0.7066} ثبت شده است.

\subsection{محدودیت‌ها و تحلیل شکست}
بررسی عملکرد ماژول بر روی معماری‌های گوناگون چالش‌های زیر را نمایان می‌سازد:
\begin{itemize}
    \item \textbf{افت کارایی در مدل خودبازگشتی \lr{GPT-2}:} ثبت دقت \lr{50.92\%} با وضعیت \lr{PASS\_LOW} روی مدل دیکودر \lr{GPT-2} نشان‌دهنده ناکافی بودن تقریب چندجمله‌ای درجه ۲ برای مهار وابستگی‌های بلندمدت و شدیداً محلی در ماسک مثلثی علّی است. تابع \lr{Softmax} با خاصیت برنده‌محور مطلق (\lr{Winner-take-all}) وزن‌های آینده را به شدت سرکوب می‌کند، در حالی که چندجمله‌ای خطی با وجود مرکزسازی، دم‌های گسترده‌تری به توکن‌های پیشین تخصیص می‌دهد که منجر به اشباع بازنمایی و توقف پیش‌بینی در آستانه حدس تصادفی شده است.
    \item \textbf{چالش‌های ادغام با مدل‌های حجیم (\lr{Qwen}):} به دلیل فرمت‌های نوین پروجکشن مشترک و عدم هماهنگی درگاه چندجمله‌ای با رفتارهای فرکانسی شدید \lr{RoPE} در بازه‌های طولانی، تزریق مستقیم ماژول بدون آموزش مقدماتی کِرنل به وضعیت گذار و لغو ارزیابی (\lr{SKIP}) منتهی شده است.
\end{itemize}
